% Standalone section fragment.  It assumes amsmath, amssymb, and amsthm,
% together with theorem, proposition, corollary, and remark environments.

\section{The all-prime Pad\'e transition and its arithmetic obstructions}
\label{sec:all-prime-pade}

Fix $0\ne\xi\in\mathbb Z$ and write
\[
 \mathcal E(u)=\sum_{m\geq0}(-1)^m m!u^m,
 \qquad
 \delta_p(\xi)=\sum_{m\geq0}(-1)^m m!\xi^m\in\mathbb Z_p.
\]
We use the classical Pad\'e pair~\cite{MZ2018}
\begin{align}
 Q_n(u)&=\sum_{k=0}^n k!\binom nk^2u^k,\label{eq:Qn-def}\\
 Q_n(u)\mathcal E(u)-P_n(u)
 &=n!^2u^{2n}\sum_{k\geq0}(-1)^k k!
       \binom{n+k}{k}^2u^k.\label{eq:pade-rem}
\end{align}
Here $Q_0=1$ and $P_0=0$.  Put
\[
 q_n(\xi)=Q_n(\xi),\qquad s_n(\xi)=P_n(\xi),\qquad
 F_{\alpha}^{(\xi)}(z)
 =\sum_{n\geq0}\bigl(q_n(\xi)\alpha-s_n(\xi)\bigr)\frac{z^n}{n!}.
\]

\begin{lemma}[Pad\'e recurrence and affine equation]
\label{lem:pade-recurrence-all-prime}
For $n\geq1$, both $q_n(\xi)$ and $s_n(\xi)$ satisfy
\begin{equation}
 u_{n+1}=\bigl(1+(2n+1)\xi\bigr)u_n-n^2\xi^2u_{n-1}.
\label{eq:pade-three-term-xi}
\end{equation}
Consequently $F=F_\alpha^{(\xi)}$ satisfies
\begin{equation}
 (1-\xi z)^2F'(z)-(1+\xi-\xi^2z)F(z)+1=0.
\label{eq:affine-pade-ode}
\end{equation}
Its normalized homogeneous solution is
\begin{equation}
 H_\xi(z)=\frac{1}{1-\xi z}
 \exp\!\left(\frac{z}{1-\xi z}\right)
 =\sum_{n\geq0}q_n(\xi)\frac{z^n}{n!}.
\label{eq:H-xi}
\end{equation}
\end{lemma}

\begin{proof}
The polynomial identity
\[
 Q_n(u)=n!u^nL_n(-1/u)
\]
and the Laguerre recurrence give
\[
 Q_{n+1}(u)=\bigl(1+(2n+1)u\bigr)Q_n(u)-n^2u^2Q_{n-1}(u).
\]
The Pad\'e determinant identity
\begin{equation}
 D_n(u):=Q_n(u)P_{n+1}(u)-Q_{n+1}(u)P_n(u)=n!^2u^{2n}
\label{eq:pade-casoratian}
\end{equation}
then yields, as a polynomial identity,
\begin{align*}
 Q_n(u)&\bigl(P_{n+1}(u)-(1+(2n+1)u)P_n(u)
                    +n^2u^2P_{n-1}(u)\bigr)\\
 &=D_n(u)-n^2u^2D_{n-1}(u)=0.
\end{align*}
Thus $P_n$ obeys the same recurrence.  Evaluation at $u=\xi$ proves
\eqref{eq:pade-three-term-xi}.

For $r_n=q_n(\xi)\alpha-s_n(\xi)$ one has
$r_0=\alpha$ and $r_1=(1+\xi)\alpha-1$.  Multiplying
\eqref{eq:pade-three-term-xi} by $z^n/n!$ and summing gives
\eqref{eq:affine-pade-ode}.  Solving its homogeneous part with value $1$
at $z=0$ gives \eqref{eq:H-xi}.
\end{proof}

\begin{theorem}[All-prime radius transition]
\label{thm:all-prime-radius}
Let $p$ be any prime, set $s=v_p(\xi)$, and let $\alpha\in\mathbb Q_p$.
Then
\begin{equation}
 R_p\!\left(F_\alpha^{(\xi)}\right)=
 \begin{cases}
 p^{\frac1{p-1}+2s}
   =p^{1/(p-1)}|\xi|_p^{-2},&\alpha=\delta_p(\xi),\\[2mm]
 p^{-1/(p-1)},&\alpha\ne\delta_p(\xi).
 \end{cases}
\label{eq:all-prime-radius}
\end{equation}
Moreover
\begin{equation}
 F_\alpha^{(\xi)}\in\mathbb Z_p[[z]]
 \quad\Longleftrightarrow\quad
 \alpha=\delta_p(\xi).
\label{eq:integrality-matching}
\end{equation}
If the right-hand condition fails, the valuations of the coefficients of
$F_\alpha^{(\xi)}$ are unbounded below; in particular no fixed nonzero
scalar makes the whole series integral.
\end{theorem}

\begin{proof}
Specializing \eqref{eq:pade-rem} at $u=\xi$ gives
\begin{equation}
 \frac{q_n(\xi)\delta_p(\xi)-s_n(\xi)}{n!}
 =n!\xi^{2n}\sum_{k\geq0}(-1)^k k!
          \binom{n+k}{k}^2\xi^k
 \in n!\xi^{2n}\mathbb Z_p.
\label{eq:matched-coefficient-divisibility}
\end{equation}
Hence the matched radius is at least
$p^{1/(p-1)+2s}$.

Let $c=1/\xi$ and $t=z-c$.  Equation
\eqref{eq:affine-pade-ode} becomes
\begin{equation}
 \xi^2t^2F'(t)+(\xi^2t-1)F(t)+1=0.
\label{eq:local-affine-ode}
\end{equation}
Its unique formal solution regular at $t=0$ is
\begin{equation}
 F_{\mathrm{reg}}(t)=\sum_{m\geq0}m!\xi^{2m}t^m,
\label{eq:local-regular-solution}
\end{equation}
whose radius is exactly $p^{1/(p-1)+2s}$.  This number is strictly larger
than $|c|_p=p^s$.  The lower bound from
\eqref{eq:matched-coefficient-divisibility} therefore places $c$ inside the
disc of the matched solution, and its re-expansion at $c$ must be
\eqref{eq:local-regular-solution}.  If its radius at $0$ were larger than
$p^{1/(p-1)+2s}$, the ultrametric discs with centers $0$ and $c$ and that
larger radius would coincide, contradicting the exact radius of
\eqref{eq:local-regular-solution}.  This proves the first case of
\eqref{eq:all-prime-radius}.

The coefficients of $H_\xi$ have radius at least $p^{-1/(p-1)}$, since
$q_n(\xi)\in\mathbb Z$.  If $p\mid\xi$, then
$q_n(\xi)\equiv1\pmod p$ for every $n$.  If $p\nmid\xi$ and
$n\equiv0\pmod p$, then the terms in \eqref{eq:Qn-def} with
$1\leq k<p$ contain $p\mid\binom nk$, while those with $k\geq p$ contain
$p\mid k!$; hence again $q_n(\xi)\equiv1\pmod p$.  Along the indicated
subsequence, Legendre's formula gives
\[
 \frac1n v_p\!\left(\frac{q_n(\xi)}{n!}\right)
 \longrightarrow-\frac1{p-1}.
\]
Thus $R_p(H_\xi)=p^{-1/(p-1)}$.  Finally,
\[
 F_\alpha^{(\xi)}-F_{\delta_p(\xi)}^{(\xi)}
   =(\alpha-\delta_p(\xi))H_\xi.
\]
When the scalar is nonzero, the term on the right has strictly smaller
radius than the matched series, proving the second case of
\eqref{eq:all-prime-radius}.  Formula
\eqref{eq:matched-coefficient-divisibility} proves integrality in the
matched case.  In the unmatched case, on the same unit subsequence,
\[
 v_p\!\left((\alpha-\delta_p(\xi))\frac{q_n(\xi)}{n!}\right)
 =v_p(\alpha-\delta_p(\xi))-v_p(n!)\longrightarrow-\infty,
\]
whereas the matched coefficient is integral.  This proves
\eqref{eq:integrality-matching} and the final assertion.
\end{proof}

\begin{proposition}[Local denominator--radius equivalence]
\label{prop:local-denominator-radius}
Put
\[
 S_n(\xi)=\sum_{k=0}^n(-1)^kk!\xi^k,
 \qquad b_n=\frac{\alpha-S_n(\xi)}{n!},
\]
and, for $\alpha\in\mathbb Q_p$, define
\[
 d_{p,N}=\max_{0\leq n\leq N}\max\{0,-v_p(b_n)\}.
\]
Then the following are equivalent:
\begin{equation}
 \alpha=\delta_p(\xi),\qquad
 b_n\in\mathbb Z_p\ (n\geq0),\qquad
 d_{p,N}=0\ (N\geq0),\qquad
 R_p(F_\alpha^{(\xi)})>1.
\label{eq:four-local-equivalences}
\end{equation}
If $\alpha\ne\delta_p(\xi)$ and
$c_p=v_p(\alpha-\delta_p(\xi))$, then, for all sufficiently large $N$,
\begin{equation}
 d_{p,N}=v_p(N!)-c_p,
 \qquad
 \lim_{N\to\infty}\frac{d_{p,N}}N=\frac1{p-1}.
\label{eq:exact-local-denominator}
\end{equation}
\end{proposition}

\begin{proof}
In the matched case,
\[
 b_n=\sum_{k\geq n+1}(-1)^k\frac{k!}{n!}\xi^k\in\mathbb Z_p.
\]
In the unmatched case, $S_n(\xi)\to\delta_p(\xi)$ and ultrametric
stability gives
$v_p(\alpha-S_n(\xi))=c_p$ for all sufficiently large $n$.  Hence
$-v_p(b_n)=v_p(n!)-c_p$ eventually.  This is nondecreasing apart from
plateaux and ultimately exceeds the maximum contributed by the finitely
many earlier terms, proving \eqref{eq:exact-local-denominator}.  The radius
equivalence is Theorem~\ref{thm:all-prime-radius}.
\end{proof}

\begin{corollary}[Relative adelic radius defect]
\label{cor:adelic-radius-defect}
Let $\alpha=a/b\in\mathbb Q$ in lowest terms and define
\[
 \mathcal B(\alpha,\xi)
 =\{p:p\nmid b\xi,\ \alpha\ne\delta_p(\xi)\}.
\]
Then
\begin{equation}
 \sum_{p\in\mathcal B(\alpha,\xi)}\frac{\log p}{p-1}=+\infty.
\label{eq:bad-prime-divergence}
\end{equation}
Consequently, with
$\rho_{p,\xi}=p^{1/(p-1)}|\xi|_p^{-2}$,
\begin{equation}
 \prod_{p\nmid b\xi}
 \frac{R_p(F_\alpha^{(\xi)})}{\rho_{p,\xi}}=0,
\label{eq:relative-adelic-product-zero}
\end{equation}
where the product is the limit over finite sets of primes.
\end{corollary}

\begin{proof}
Suppose that the sum in \eqref{eq:bad-prime-divergence} were finite, and let
$\mathcal P$ be the primes $p\nmid b\xi$ for which
$\delta_p(\xi)=\alpha$.  The complement of $\mathcal P$ then has finite
$\sum_p\log p/(p-1)$.  Since
\[
 \prod_{p\in\mathcal P}|n!|_p^2
 =\frac1{(n!)^2}\prod_{p\notin\mathcal P}p^{2v_p(n!)}
 \leq\frac{C^n}{(n!)^2}
\]
for a fixed $C>0$, the hypothesis of the global-relation theorem of
Matala-aho and Zudilin~\cite{MZ2018} is satisfied.  That
theorem forbids all values $\delta_p(\xi)$, $p\in\mathcal P$, from being the
same rational number, a contradiction.  Thus
\eqref{eq:bad-prime-divergence} holds.  For $p\nmid b\xi$, Theorem
\ref{thm:all-prime-radius} gives a local ratio $1$ at a matched prime and
$p^{-2/(p-1)}$ at an unmatched prime.  Taking logarithms proves
\eqref{eq:relative-adelic-product-zero}.
\end{proof}

\begin{remark}[What the adelic product does and does not say]
The product in \eqref{eq:relative-adelic-product-zero} is normalized by the
matched local benchmark.  It is not itself an application of the
Borel--Dwork--P\'olya--Bertrandias rationality criterion.  Rather, it records
why the direct application fails here: for rational $\alpha$ there are
infinitely many bad integral lattices, with divergent total radius defect.
It does not exclude a different auxiliary series or different adelic
domains; compare the original Dwork criterion~\cite{Dwork1960} and its
capacity-theoretic extension~\cite{BostChambertLoir2009}.
\end{remark}

\subsection{Frobenius and Borel obstructions}

\begin{proposition}[Exact $p$-curvature]
\label{prop:exact-p-curvature}
Let $u=1-\xi z$.  The homogeneous equation in
Lemma~\ref{lem:pade-recurrence-all-prime} is the rank-one connection
\[
 \nabla_\partial=\partial-a_\xi(z),
 \qquad
 a_\xi(z)=\frac{\xi}{u}+\frac1{u^2}.
\]
At every prime $p$, its reduction has $p$-curvature
\begin{equation}
 \psi_p(\partial)=-\frac1{(1-\bar\xi z)^{2p}},
\label{eq:exact-p-curvature}
\end{equation}
up to the conventional overall sign.  In particular it is nonzero, hence
not nilpotent, at every prime.
\end{proposition}

\begin{proof}
In characteristic $p$, the rank-one Jacobson formula gives
\[
 (\partial-a_\xi)^p
 =-\bigl(\partial^{p-1}a_\xi+a_\xi^p\bigr),
\]
because $\partial^{[p]}=0$ on $\mathbb F_p(z)$.  Wilson's congruence gives
\[
 \partial^{p-1}(\xi u^{-1})=-\xi^pu^{-p},
 \qquad
 \partial^{p-1}(u^{-2})=p!\xi^{p-1}u^{-p-1}=0,
\]
whereas Frobenius gives
\[
 a_\xi^p=\xi^pu^{-p}+u^{-2p}.
\]
Substitution proves \eqref{eq:exact-p-curvature}.
\end{proof}

\begin{corollary}[Global-nilpotence no-go]
\label{cor:global-nilpotence-no-go}
The minimal rank-one operator annihilating $H_\xi$ is not globally
nilpotent.  Consequently it is not a subquotient of an ordinary smooth
proper Gauss--Manin connection and is not a $G$-operator in the usual
globally nilpotent category.
\end{corollary}

\begin{proof}
For a rank-one connection, nilpotence of $p$-curvature is equivalent to its
vanishing.  Proposition~\ref{prop:exact-p-curvature} therefore rules out
global nilpotence.  The stated consequences are the standard
Gauss--Manin/global-nilpotence theorems of Katz~\cite{Katz1970} and the
Andr\'e--Chudnovsky--Katz theory of $G$-operators~\cite{Andre1989}.
\end{proof}

\begin{remark}
This is deliberately a statement about the ordinary globally nilpotent
category.  It does not rule out every notion of irregular exponential
geometric origin.  Moreover the connection is independent of $\alpha$:
$p$-curvature cannot distinguish the matched initial value.  Matching is a
property of an overconvergent horizontal lift, not of the differential
module alone.
\end{remark}

\begin{proposition}[Dieudonn\'e--Dwork defect]
\label{prop:dwork-defect}
Put $w(z)=z/(1-\xi z)$.  For every prime $p$,
\begin{equation}
 [z^p]\bigl(w(z^p)-pw(z)\bigr)=1-p\xi^{p-1}\notin p\mathbb Z_p.
\label{eq:dwork-defect-coefficient}
\end{equation}
Consequently neither $\exp(w(z))$ nor $H_\xi(z)$ belongs to
$\mathbb Z_p[[z]]$.
\end{proposition}

\begin{proof}
The Dieudonn\'e--Dwork criterion~\cite{Dieudonne1957} says that, for
$S\in z\mathbb Q_p[[z]]$,
\[
 \exp(S)\in1+z\mathbb Z_p[[z]]
 \quad\Longleftrightarrow\quad
 S(z^p)-pS(z)\in pz\mathbb Z_p[[z]].
\]
Since $w(z)=\sum_{m\geq1}\xi^{m-1}z^m$, equation
\eqref{eq:dwork-defect-coefficient} follows immediately and violates the
criterion.  Finally $(1-\xi z)^{-1}$ and its inverse both lie in
$\mathbb Z_p[[z]]$, so multiplication by this factor cannot repair
integrality.
\end{proof}

\begin{proposition}[Fixed Borel logarithm]
\label{prop:fixed-borel-logarithm}
Let
\[
 \frac{E_\xi(z)-\alpha}{z-1}=\sum_{n\geq0}n!b_nz^n,
 \qquad
 E_\xi(z)=\sum_{n\geq0}(-1)^nn!\xi^nz^n,
\]
and let $\mathcal B_1$ denote the order-one formal Borel transform.  Then
\[
 B_{\alpha,\xi}(t):=
 \mathcal B_1\!\left(\frac{E_\xi-\alpha}{z-1}\right)(t)
 =\sum_{n\geq0}b_nt^n
\]
satisfies
\begin{equation}
 B_{\alpha,\xi}'-B_{\alpha,\xi}
 =\frac{\xi}{(1+\xi t)^2},
 \qquad B_{\alpha,\xi}(0)=\alpha-1.
\label{eq:borel-affine-ode}
\end{equation}
If $t_0=-1/\xi$ and $x=t-t_0$, then on every local complex branch
\begin{equation}
 B_{\alpha,\xi}(t)
 =-\frac{e^x}{\xi x}-\frac{e^x}{\xi}\log x+e^xh_{\alpha,\xi}(x),
 \qquad h_{\alpha,\xi}\in\mathbb C\{x\}.
\label{eq:fixed-borel-log}
\end{equation}
In particular no choice of $\alpha$ makes this Borel transform meromorphic
at $t_0$, and none makes it rational.
\end{proposition}

\begin{proof}
Coefficient comparison gives
$b_n=(\alpha-S_n(\xi))/n!$ and hence
\[
 (n+1)b_{n+1}=b_n+(n+1)(-1)^n\xi^{n+1}.
\]
Summation proves \eqref{eq:borel-affine-ode}, equivalently
\[
 B_{\alpha,\xi}(t)=e^t\left(
 \alpha-1+\int_0^t\frac{\xi e^{-v}}{(1+\xi v)^2}\,dv\right).
\]
At $t=t_0+x$, equation \eqref{eq:borel-affine-ode} reads
$B'-B=1/(\xi x^2)$.  Equivalently,
\[
 (e^{-x}B)'=\frac{e^{-x}}{\xi x^2}.
\]
Since
$e^{-x}x^{-2}=x^{-2}-x^{-1}+\text{a holomorphic germ}$, integration and
multiplication by $e^x$ give \eqref{eq:fixed-borel-log}.  The homogeneous
correction $Ce^t$, which is the only dependence on $\alpha$, is holomorphic
and cannot cancel either singular term.
\end{proof}

\begin{corollary}[Polynomial reduction in the matched case]
\label{cor:matched-polynomial-reduction}
If $p\nmid\xi$, then
\begin{equation}
 \overline{F_{\delta_p(\xi)}^{(\xi)}}(z)
 =\sum_{m=0}^{p-1}m!\bar\xi^{2m}
       (z-\bar\xi^{-1})^m\in\mathbb F_p[z].
\label{eq:matched-polynomial-mod-p}
\end{equation}
If $p\mid\xi$, then
$\overline{F_{\delta_p(\xi)}^{(\xi)}}(z)=1$.  Thus the matched reduction
has a finite-support, hence $p$-automatic, coefficient sequence.  In the
unmatched case no fixed scalar gives an integral series to reduce.
\end{corollary}

\begin{proof}
For $p\nmid\xi$, reduce the local expansion
\eqref{eq:local-regular-solution}; all terms with $m\geq p$ vanish modulo
$p$, giving \eqref{eq:matched-polynomial-mod-p}.  If $p\mid\xi$, equation
\eqref{eq:matched-coefficient-divisibility} shows that every positive-degree
coefficient is divisible by $p$, while
$F_{\delta_p(\xi)}^{(\xi)}(0)=\delta_p(\xi)\equiv1\pmod p$.  The last two
claims follow from Theorem~\ref{thm:all-prime-radius}; the automaticity
statement is also the elementary polynomial case of Christol's
theorem~\cite{ChristolEtAl1980}.
\end{proof}

\begin{remark}[Kurepa specialization]
For $\xi=-1$,
\[
 \delta_p(-1)\equiv\sum_{m=0}^{p-1}m!=!p\pmod p,
 \qquad
 \overline{F_{\delta_p(-1)}^{(-1)}}(z)
 =\sum_{m=0}^{p-1}m!(z+1)^m.
\]
Thus Kurepa's conjecture is equivalent to the assertion that the constant
term of this matched polynomial is nonzero for every odd prime.  This is an
exact reformulation, not a proof of the conjecture.
\end{remark}

% Bibliography keys required by this fragment:
% MatalaAhoZudilin2018 -- T. Matala-aho and W. Zudilin,
%   Euler's factorial series and global relations, JNT 186 (2018), 202--210.
% Dieudonne1957 -- J. Dieudonne, On the Artin--Hasse exponential series,
%   Proc. AMS 8 (1957), 210--214.
% Katz1970 -- N. Katz, Nilpotent connections and the monodromy theorem,
%   Publ. Math. IHES 39 (1970), 175--232.
% ChristolEtAl1980 -- G. Christol, T. Kamae, M. Mendes France, G. Rauzy,
%   Suites algebriques, automates et substitutions, BSMF 108 (1980), 401--419.
% Dwork1960 -- B. Dwork, On the rationality of the zeta function of an
%   algebraic variety, Amer. J. Math. 82 (1960), 631--648.
% BostChambertLoir2009 -- J.-B. Bost and A. Chambert-Loir, Analytic curves
%   in algebraic varieties over number fields, in Algebra, Arithmetic, and
%   Geometry, Progress in Mathematics 269 (2009), 69--124.
% Andre1989 -- Y. Andre, G-functions and Geometry, Aspects of Mathematics
%   E13, Vieweg, 1989.
